% Minimal preamble for P2 cards taken from sources/comparability-partitions-2408.13523v2/source0.tex
% (arXiv:2408.13523v2, "Partitioning perfect graphs into comparability graphs").
% Every line below is copied mechanically from that file; nothing was invented, renamed or redefined.
%   line 2    : \documentclass[12pt]{article}
%   line 3    : \usepackage{amssymb,latexsym}
%   line 4    : \usepackage{amsmath}
%   line 5    : \usepackage[dvipdfm]{graphicx}  -- the dvipdfm driver option is dropped, because the
%               assembler forces the xelatex engine and a forced driver option contradicts it.
%   line 16   : \usepackage{amsfonts}
%   line 17   : \usepackage{amsthm}
%   lines 20-30: the author's theorem declarations
%   lines 32-41: the author's own macro definitions
% Omitted source preamble lines and why they are not needed for these body spans:
%   line 1 (provenance comment), lines 6-11 (tikz and its libraries: no drawing in the retained spans),
%   lines 7 and 9 (verbatim), lines 12-13 (page geometry), line 14 (rawfonts),
%   line 15 (psfrag), line 18 (hyperref with the pagebackref option: the retained spans print no
%   citation anchors and the assembly adds its own hyperref machinery), lines 43-61 (the author's
%   backref neutralisation and his margin-note macros, which depend on the omitted hyperref).
% The retained spans of the body use only \em, \bf, \cal, amsthm environments, amsmath displays and
% the macros defined in lines 32-41 below.
\documentclass[12pt]{article}
\usepackage{amssymb,latexsym}
\usepackage{amsmath}
\usepackage{graphicx}
\usepackage{amsfonts}
\usepackage{amsthm}

\newtheorem{theorem}{Theorem}%[section]
\newtheorem{problem}[theorem]{Problem}%[section]
\newtheorem{observation}[theorem]{Observation}
\newtheorem{lemma}[theorem]{Lemma}
\newtheorem{corollary}[theorem]{Corollary}
\newtheorem{definition}{Definition}
\newtheorem{example}[theorem]{Example}
\newtheorem{conjecture}{Conjecture}%[section]
\newtheorem{proposition}[theorem]{Proposition}
\newtheorem{claim}[theorem]{Claim}
\newtheorem{remark}[theorem]{Remark}

\def\qed{\ifhmode\unskip\nobreak\fi\quad\ifmmode\Box\else$\Box$\fi}
\def\pont{\hspace{-6pt}{\bf.\ }}
\def\qed{\hfill $\Box$}

\def\R{{\mbox{\boldmath$\mathcal R$}}}
% the author's second definition of \R, left commented out exactly as in the source
\def\F{{\mathcal F}}
\def\P{{\mathcal P}}

\def\newpage{\vfill\eject}
\title{Every unit interval graph has its edge set partitioned into two comparability graphs}
\author{András Gyárfás, Márton Marits, Géza Tóth}
\date{arxiv:2408.13523v2; CC BY 4.0}
\begin{document}
\maketitle
\noindent\textit{Editorial note.} The counted unit is the article's own Theorem unit at source lines 281-282 together with the author's proof that opens with the bold header Proof. at source line 285 and closes with the square at source line 300; the statement environment ends at line 282 and only blank lines and a medskip lie between, so the pairing is the ordinary adjacent one.\par\medskip
\noindent\textit{Editorial note.} The retained context at source lines 117-120 and 128-129 is the article's own description of a comparability graph and its own definitions of the parameters p(G) and c(G); the retained context at source lines 272-273 is its own definition of interval graphs and of unit interval graphs.\par\medskip
\noindent\textit{Editorial note.} The source preamble cannot be retained as it stands: it loads rawfonts, psfrag, verbatim, tikz with several libraries, hyperref with the pagebackref option and a dvi driver option for graphicx, none of which the retained body spans use, and it redefines backref commands that only exist because of that hyperref option. The retained preamble file therefore carries, line for line, source line 2, source line 3, source line 4, source line 16, source line 17, the theorem declarations of source lines 20-30 and the author macro definitions of source lines 32-41; the only textual change anywhere is the removal of the dvipdfm driver option from the graphicx line, because the assembly forces the xelatex engine. No new macro, package or declaration was introduced.\par\medskip
\noindent\textit{Editorial note.} No retained line contains a citation or a cross-reference: the sentences of the article that cite Gilmore and Hoffman, Gavril and Walter are deliberately not retained, so the material carries no bibliography and needs no external navigation pointer.\par\medskip
\noindent\textit{Editorial note.} The article's theorems co-triangulated, gensplit, bipline, loglog and ub, its lemma rec and its proposition H9 are not counted here and are not used as premises.\par\medskip
\noindent\textit{Allowed standard prerequisites (editorial).}\begin{itemize}
\item Comparability graph of a partially ordered set. Given a partially ordered set, its comparability graph has the elements as vertices and an undirected edge between two vertices exactly when the corresponding elements are comparable; the article's own description is retained at source lines 117-120.. Fixes the class of graphs into which the proof must partition the edge set of the interval graph, namely two edge-disjoint comparability subgraphs.
\item The partition number p(G). p(G) is the minimum m such that the edge set of G is the union of m edge-disjoint comparability subgraphs of G; retained at source line 128, together with the covering number c(G) at source line 129.. The counted theorem asserts p(G) <= 2, so the parameter and its edge-disjointness requirement must be fixed before the statement can be read.
\item Interval graph and unit interval graph. A graph is an interval graph when it is the intersection graph of a system of closed intervals of a line, and a unit interval graph when the intervals all have length one; the article's own definitions are retained at source lines 272-273.. Identifies the graph G of the counted theorem with the intersection graph G\_U of a system U of unit intervals, which is the object the proof partitions.
\end{itemize}\medskip
Here we focus on a well known family of perfect graphs, the {\em comparability graphs},
which can be obtained from a partially ordered set as follows.
Vertices correspond to the elements of the partially ordered set,
and two vertices are connected by an undirected edge if and only if the corresponding elements

Let $p(G)$ denote the minimum number $m$ such that $E(G)=\cup_{i=1}^m E(G_i)$ where the $G_i$-s are edge-disjoint comparability subgraphs of $G$.
Dropping  the condition that the $G_i$-s are edge-disjoint, we define $c(G)$ as the smallest number of comparability subgraphs of $G$ needed to {\em cover} $E(G)$.

A graph $G$ is an {\em interval graph} if it is the intersection graph of a system of closed intervals of a line.
If the intervals have length one, their intersection graph is a {\em unit interval graph}.

\begin{theorem}\label{unit} For unit interval graphs $G$, $p(G)\le 2$.
\end{theorem}

\noindent{\bf Proof. }  Let $U$ be a system of unit intervals and partition them as follows.
Let $I_1=[a_1,b_1] \in U$ be the interval with leftmost right endpoint and define $X_1$ as the set of intervals of
$U$ containing $b_1$.
If $I_1,\dots, I_k$ and $X_1,\dots,X_k$ are defined and there are further intervals, select from them the interval $I_{k+1}=[a_{k+1},b_{k+1}]$
with the leftmost right endpoint.
Define $X_{k+1}$ as the set of intervals not in $\cup_{i=1}^k X_i$ and containing $b_{k+1}$.
The procedure ends with $[a_m,b_m]$ and $X_m$ when all intervals of $U$ are partitioned.

The unit interval graph of $U$, $G_U$, can be partitioned into two comparability graphs as follows.
Define $H_1$ as the edge set of $G_U$ defined by intervals of the {\em same}
$X_i$. Note that $H_1$ is a union of $m$ vertex-disjoint complete graphs.
Let $H_2$ be the set of remaining edges, i.e. whose intersecting intervals belong to {\em different}
$X_i$-s.
Because $U$ contains unit intervals,
edges of $H_2$ must represent intersecting intervals belonging to $X_i$ and $X_{i+1}$ for some $i$.
Thus $H_2$ is a bipartite graph. Clearly $H_1,H_2$ are both comparability graphs, finishing the proof.   \qed

\end{document}
